\section{A deterministic support certificate}
\label{sec:tpc53-support}

The dictionary is exact even when the physical profile is zero.  A
lower conditioning bound requires a separate overlap statement.  We now
make that statement finite-dimensional and checkable.

\subsection{Continuous carrier and profile energies}

Let \(Y=[0,\log2]^2\), with interior \(Y^\circ=(0,\log2)^2\), and
\begin{equation}
 d\nu(q)=\e^{x+y}\,dx\,dy,
 \qquad q=(x,y).
 \label{eq:tpc53-row-limit-measure}
\end{equation}
This measure has mass one on \(Y\).  For a compact source-parameter set
\(\mathfrak B\), equipped with its relative topology in the normalized
parameter domain, define
\begin{align}
 \mathscr B(\beta)
 &:=\frac1{H_0}\int_Y\int_{I_0}
  |B_{\beta,a}(q,t)|^2\,dt\,d\nu(q),
 \label{eq:tpc53-continuous-carrier-energy}\\
 \mathscr P(\beta)
 &:=\frac1{H_0}\int_Y\int_{I_0}
  |F_{\beta,a}(q,t)|^2\,dt\,d\nu(q).
 \label{eq:tpc53-continuous-profile-energy}
\end{align}
We discard a zero carrier cell.  Compactness then gives
\begin{equation}
 0<b_-\le\mathscr B(\beta)\le b_+<\infty
 \quad(\beta\in\mathfrak B)
 \label{eq:tpc53-carrier-bounds}
\end{equation}
after restricting to a nonzero compact source cell.

For a smooth function \(U\), write
\begin{equation}
 S_U:=\{z>0:U(z)\ne0\}.
 \label{eq:tpc53-nonzero-set}
\end{equation}
Every \(S_U\) is open.

\begin{definition}[Strictly reachable point]
\label{def:tpc53-strict-point}
A point \((\beta_0,q_0,t_0)\in
\mathfrak B\times Y^\circ\times\operatorname{int}I_0\)
is strictly reachable when
\begin{equation}
 B_{\beta_0,a}(q_0,t_0)\ne0,
 \qquad
 z_1\in S_{W_1},\quad z_2\in S_{\Psi_1},\quad
 z_3\in S_{W_2},\quad z_4\in S_{\Psi_2},
 \label{eq:tpc53-strict-reachability}
\end{equation}
with every argument evaluated at \((\beta_0,q_0,t_0)\).
\end{definition}

For the two opposite-row arguments, strict reachability implies the
useful ratio condition
\begin{equation}
 \frac{z_3}{z_4}=\frac{c_X}{c'_X}\e^{x_0}
 \in\left(\frac{c_X}{c'_X},
          2\frac{c_X}{c'_X}\right).
 \label{eq:tpc53-support-ratio}
\end{equation}
At the TPC--25 scaling \(c'_X=1\), this is a direct comparison between
the nonzero sets of \(W_2\) and \(\Psi_2\).

\begin{theorem}[Reachable-support equivalence]
\label{thm:tpc53-support-equivalence}
For a fixed \(\beta\) satisfying
\eqref{eq:tpc53-carrier-bounds}, the following are equivalent:
\begin{enumerate}[label=\textup{(\roman*)},leftmargin=3em]
 \item \(\mathscr P(\beta)>0\);
 \item there is a strictly reachable point with source parameter
 \(\beta\);
 \item there are compact boxes \(\cQ\Subset Y^\circ\),
 \(I\Subset\operatorname{int}I_0\), and constants
 \(r_0,w_{1,0},\psi_{1,0},w_{2,0},\psi_{2,0}>0\) such that throughout
 \(\cQ\times I\),
 \begin{equation}
 |B_{\beta,a}|\ge r_0,\quad
 |W_i(z_{2i-1})|\ge w_{i,0},\quad
 |\Psi_i(z_{2i})|\ge\psi_{i,0}\quad(i=1,2).
 \label{eq:tpc53-box-lower-bounds}
 \end{equation}
\end{enumerate}
\end{theorem}

\begin{proof}
The integrand in \eqref{eq:tpc53-continuous-profile-energy} is
continuous and nonnegative, and the measure
\(d\nu(q)dt\) is positive on every nonempty interior box.  A positive
integral therefore has a point where every factor is nonzero, proving
\textup{(i)} implies \textup{(ii)}.  Openness of all nonzero sets and
continuity give a compact product neighborhood and positive minima,
which is \textup{(iii)}.  Integrating the positive lower bound in
\textup{(iii)} proves \textup{(i)}.
\end{proof}

\begin{theorem}[Explicit compact-box overlap]
\label{thm:tpc53-explicit-overlap}
Suppose one strict point exists with
\(\beta_0\in\operatorname{relint}\mathfrak B\).  Then there are
\(\mathfrak B_0\Subset\mathfrak B\), \(\cQ\Subset Y^\circ\), and
\(I\Subset\operatorname{int}I_0\) on which
\eqref{eq:tpc53-box-lower-bounds} holds uniformly in
\(\beta\in\mathfrak B_0\).  For the full continuous energies, with the
source parameter restricted to \(\mathfrak B_0\),
\begin{equation}
 \boxed{
 \rho_*:=\inf_{\beta\in\mathfrak B_0}
 \frac{\mathscr P(\beta)}{\mathscr B(\beta)}
 \ge
 \frac{r_0^2w_{1,0}^2\psi_{1,0}^2w_{2,0}^2\psi_{2,0}^2}
      {H_0b_+}
 \left(\int_{\cQ}\e^{x+y}\,dx\,dy\right)|I|>0.}
 \label{eq:tpc53-explicit-rho}
\end{equation}
\end{theorem}

\begin{proof}
Continuity in \(\beta\) extends the strict inequalities to a compact
source box.  The numerator of the profile energy is at least the
integral over \(\cQ\times I\) of the product of the squared lower bounds.
The carrier denominator is at most \(b_+\).  This gives
\eqref{eq:tpc53-explicit-rho}.
\end{proof}

Thus the active source set
\begin{equation}
 \mathfrak B_{\rm act}
 :=\{\beta:\mathscr P(\beta)>0\}
 \label{eq:tpc53-active-set}
\end{equation}
is relatively open.  Every compact subset of it is uniformly active.  A sequence
of active parameters whose reachable arguments graze a cutoff boundary
can lose its lower bound at the boundary; compactness of the full closed
scale range alone therefore gives no positive uniform bound.

\subsection{A legal nonempty model cell}

The certificate is not vacuous.  Take four nonnegative smooth cutoffs
supported in \((1/2,2)\) and equal to one on \([3/4,5/4]\), take the
fixed opposite unary shape \(R\) to equal one near the row point below,
and take \(c_X=c'_X=1\).  Choose \(I_0\) to contain the point below.  With
\[
 x_m=y_m=x_\gamma=y_\gamma=\log(21/20),
 \qquad t_0=(21/20)^{-2},
\]
the two \(W\)-arguments equal one and the two \(\Psi\)-arguments equal
\(20/21\).  This is a strict interior point.  It proves that the allowed
smooth cutoff class contains certified packets; it does not assert that
an unspecified inherited cutoff has these supports.
